\begin{lemma}
\label{lemma-irreducible-polynomials}
Let $F$ be a field. Let $P \in F[x]$ be an irreducible polynomial over $F$.
Let $P' = \text{d}P/\text{d}x$ be the derivative of $P$ with respect
to $x$. Then one of the following two cases happens
\begin{enumerate}
\item $P$ and $P'$ are relatively prime, or
\item $P'$ is the zero polynomial.
\end{enumerate}
The second case can only happen if $F$ has characteristic $p > 0$.
In this case $P(x) = Q(x^q)$ where $q = p^f$ is a power of $p$ and
$Q \in F[x]$ is an irreducible polynomial such that $Q$ and $Q'$
are relatively prime.
\end{lemma}

\begin{proof}
If $P' = 0$, we are in case (2). Otherwise
$\deg(P') < \deg(P)$. If $P$ and $P'$ were not relatively prime,
write $(P, P') = (R)$. Then $R$ is a nonunit divisor of $P$,
so irreducibility gives $(R) = (P)$. But $R$ also divides the
nonzero polynomial $P'$, which would force
$\deg(P) = \deg(R) \leq \deg(P') < \deg(P)$, a contradiction.
Thus $P$ and $P'$ are relatively prime, and we are in case (1).

\medskip\noindent
Assume we are in case (2) and $P = a_d x^d + \ldots + a_0$. Then
$P' = da_d x^{d - 1} + \ldots + a_1$. In characteristic $0$ we see
that this forces $a_d, \ldots, a_1 = 0$ which would mean $P$ is constant
a contradiction. Thus we conclude that the characteristic $p$ is positive.
In this case the condition $P' = 0$ forces $a_i = 0$ whenever $p$ does
not divide $i$.
In other words, $P(x) = P_1(x^p)$ for some nonconstant polynomial $P_1$.
Clearly, $P_1$ is irreducible as well. By induction on the degree we
see that $P_1(x) = Q(x^q)$ as in the statement of the lemma, hence
$P(x) = Q(x^{pq})$ and the lemma is proved.
\end{proof}